% =====================================================================
%  c01.tex —— §1 引言（Lead）
%  原文 PDF pp.3--8（印刷 pp.2--7）
% =====================================================================
% ===== 原文 PDF pp.3--8（印刷 pp.2--7）=====

% [原文 p.2]
\SEC{1}{引言}

约一个世纪之前诞生以来，四维时空中的量子场论（QFT）一直是理论物理中一个重要而持久的主题。
时至今日，它已经给出了极其精确的预言，令人惊叹地准确描述了我们的物理世界，
例如大型强子对撞机（LHC）上的粒子碰撞。然而，我们对 QFT 在根本层面上的理解仍然相当有限。

QFT 的常规进路依赖于「基本量子场 + 刻画相互作用的拉格朗日量」这一表述；在此框架下，
关联函数、散射振幅等物理可观测量可以通过微扰费曼图计算得到。但人们很快意识到，
对于一般的 QFT，这类微扰方法要么失效、要么根本不存在：当系统处于强耦合、
且相关物理可观测量没有（显然的）小展开参数时，即属此列。在这种情形下，
非微扰效应变得重要，而拉格朗日量方法已不再有用。由于规范理论的渐近自由
——包括描述束缚夸克的强相互作用的量子色动力学（QCD）在内——这类强耦合现象
在四维时空 QFT 的红外极限中尤其常见。显而易见的挑战在于：发展出能够辅助并超越
拉格朗日量方法的非微扰手段。这已成为近年来许多研究工作的焦点，并在四维中催生出
格外丰富多样的技术。

驯服 QFT 中强耦合动力学的一大抓手来自超对称（SUSY）。与任何对称性一样，
SUSY 能澄清理论的相图，并约束那些保持（部分）SUSY 的可观测量；其效率之高，
使得理解强耦合区域的物理变得容易得多。一个典型例子是文献 [1--4] 中对
$\cN=1$ 超 QCD（SQCD）类理论的研究。这里 $\cN$ 计量理论中超对称的个数；
人们自然预期，$\cN$ 越大，来自 SUSY 的约束越强。对于四维的相互作用 QFT，
$\cN$ 只能取 $1$ 到 $4$ 之间的值 [5]。在极端的 $\cN=4$ 情形，人们相信理论由
规范群唯一确定（至多相差离散拓扑数据），并对应于 $\cN=4$ 超杨--米尔斯理论。%
\footnote{特别地，$\cN=4$ 超共形场论（SCFT）都带有一个复的恰当边缘耦合 [5,6]，
它与超杨--米尔斯（SYM）理论中的复化规范耦合
$\tau=\frac{4\pi\ii}{g_{\mathrm{YM}}^2}+\frac{\theta}{2\pi}$ 等同。
对固定的规范群，SYM 还可能带一个额外的离散 theta 角 [7]。
我们注意到，「所有 $\cN=4$ SCFT 都是 $\cN=4$ SYM」的严格证明至今仍未给出；
利用超共形自举已取得了一些进展 [8,9]。}
% [Tutor Note: 原文 "N can take values between 1 and 4 [5]" 指四维相互作用理论的超对称数
%  上限为 4：$\cN>4$ 的（无质量、自旋 $\le 1$）超多重态必含自旋 $>1$ 的场，与可重整
%  相互作用理论不符。$\cN=3$ 的理论与 $\cN=4$ 在其超共形代数层面等价（$\cN=3$
%  超共形代数在特定条件下自动增强为 $\cN=4$），故真正独立的极大超对称情形只有 $\cN=4$。]
% [原文 p.3]
% [Tutor Note: 上一条脚注给出的 $\tau=\frac{4\pi\ii}{g_{\mathrm{YM}}^2}+\frac{\theta}{2\pi}$
%  是 SYM 的复化规范耦合，其中 $\theta$ 项为拓扑项 $\frac{\theta}{8\pi^2}\int\Tr(F\wedge F)$。
%  对 $\cN=4$ 而言 $\beta$ 函数恒为零，故 $\tau$ 是恰当边缘的——这正是原文说
%  「all have a complex exactly marginal coupling」的意思；任意 $\tau$ 的上半平面都属于同一
%  共形流形，模去 $\SL(2,\Z)$ 的离散identification后得到基本区域。]

$\cN=4$ SYM 在 70 年代末至 80 年代初被大量研究，特别是为了证明它是紫外有限的、
且共形不变 [10--13]。%
\footnote{我们感谢 Peter West 向我们解释了四维 SCFT 的早期历史。}
事实上，$\cN=4$ SYM 是四维中发现的第一个共形场论（CFT），
此后不久，人们便探索了向超对称较少的杨--米尔斯理论的推广（例如见 [14--19] 中
关于有限性的类似分析）。正如 [20,21] 的工作所表明的那样，$\cN=2$ 的情形恰好位于
「SUSY 的强大约束」与「四维中令人感兴趣的强耦合动力学」之间的一块甜蜜点。
一方面，超出 $\cN=1$ 的那部分额外 SUSY 提供了对非微扰效应（如规范理论中的瞬子与磁单极子）
进行定量控制所必需的手段，从而得以在强耦合处提取物理可观测量；另一方面，
$\cN=2$ QFT 与四维中更一般的强耦合理论共享许多特征，例如涌现的共形对称性与电磁对偶 [22,23]。
% [Tutor Note: 「紫外有限」指 $\cN=4$ SYM 的 $\beta$ 函数在所有阶都为零，
%  这源于其 $\cN=4$ 矢量多重态恰好由 $\cN=1$ 矢量多重态 $\oplus$ 三个 $\cN=1$ 手征多重态组成，
%  后者贡献的物质部分刚好抵消规范部分的 $\beta$ 函数。历史上的证明用到了
%  「$\cN=4$ 是 $\cN=1$ 的 $\beta$ 函数的零点」以及实例计算 [10--13]。]

尽管四维 $\cN=2$ QFT 最初是基于拉格朗日量方法 [15,16,20,21,24] 被构造和研究的，
但「一般的四维 $\cN=2$ QFT 究竟是什么」这一图像多年来有了飞跃式的发展：
这得益于 [22,23,25] 中强耦合理论的场论构造，以及随后来自弦/M/F 理论中
各种各样的推广 [26--60]。考虑存在保持四维 $\cN=2$ Poincaré 超对称的膜与通量的奇点几何，
在引力退耦的某些极限下（例如弦标度与弦耦合的极限），有意思的场论便从中涌现出来。
虽然其中一些理论具有熟悉的 $\cN=2$ 拉格朗日量，但按此方式产生的一般 $\cN=2$ 理论
是非拉格朗日型的；尽管如此，凭借弦对偶 [61]，这些理论（包括所有非微扰效应在内）
都被弦/M/F 理论构造完全确定。因此，场论的强耦合相往往可以直接从几何中读出。
特别地，描述 $\cN=2$ 超对称不动点的 $\cN=2$ SCFT，来自弦/M/F 理论中几何构造的某些极限；
这些极限具有标度对称性，而该对称性在所得场论中成为完整超共形对称性的一部分。
既然 SCFT 及其形变勾勒出了四维 $\cN=2$ 理论的图景，我们就有必要去探索这些不动点
在弦/M/F 理论几何构造中的一般实现。
% [Tutor Note: 「非拉格朗日型」并不等于「不可定义」：这类理论没有已知的、以基本场写出的
%  弱耦合拉格朗日量，但其完整的算子谱与非微扰数据由几何（奇点 + 膜 + 通量）以及
%  弦对偶唯一确定。理解这一点是后文 §4（Class S，来自 M5 膜）与 §5（几何工程，来自 IIB 型
%  探针三重复形奇点）两条主线的共同出发点。]
% [原文 p.4]

这里我们来到这组讲义的主要目的：即简要介绍弦/M/F 理论中一组经过挑选的几何工具，
这些工具既能构造四维 $\cN=2$ SCFT，又能确定其中的物理可观测量。
一种进路依赖于 [33,35] 中引入的 Class S 构造的推广版本：它使用 M5 膜
（更准确地说，是其世界体上的六维 $\cN=(2,0)$ 理论），紧化在带穿刺点与扭曲的黎曼面上，
我们将在 §4 中回顾。另一种进路则涉及用 IIB 型弦论去探针 [31] 首先引入的孤立三重复形奇点，
在 §5 中回顾。我们还会在 §5.4.2 中简要讨论 F 理论构造，它允许非平凡的轴子--伸缩子背景，
因而是对 IIB 型设定的推广。每种构造都给出一族无穷多的四维 $\cN=2$ SCFT；
而同一个 SCFT 常常可以从多种弦/M/F 理论构造中得到，这些构造往往为同一个 SCFT
提供互补的视角。这样挑选主题，是为了对近期若干关于四维 $\cN=2$ CFT 的综述 [62--64]
作一点得体的补充。我们强调，另有一些四维 $\cN=2$ SCFT 的构造不在本文的讨论范围之内，
例如来自六维 $\cN=(1,0)$ 或五维 $\cN=1$ SCFT 的环面紧化，代表性文献见 [65--71]。%
\footnote{特别地，系统地理解四维 $\cN=2$ SCFT 的「直接几何工程构造」与
「由六维、五维 SCFT 紧化得到的构造」之间的关系将非常有意思；后者在弦/M/F 理论中
也各有其几何构造。}
最后我们将在 §6 中讨论若干开放问题。

为继续展开，我们先在 §2 中回顾四维 $\cN=2$ 超对称与超共形对称性的基础，
并关注 $\cN=2$ SCFT 中重要的物理可观测量，例如受保护的算子谱与反常。
由于缺少微扰拉格朗日量描述，这些可观测量在不动点处通常很难直接得到。
一个自然的策略是研究不动点理论的普适形变，并通过外推形变理论中的可观测量来提取 SCFT 数据。
$\cN=2$ SCFT 的一个普遍特征是：存在一个保持 $\cN=2$ SUSY 的真空模空间，
从而可以与拉格朗日量理论建立有用的联系。特别地，我们考虑这样的 SCFT 形变——
即移动到模空间中所谓的库仑分支（CB）上；对于相互作用的 $\cN=2$ SCFT，这样的分支一般被期待存在。
库仑分支上的远红外（IR）物理由一个 $\cN=2$ 阿贝尔规范理论描述，其相互作用由全纯预势支配
——这是 $\cN=2$ SUSY 的结果，我们将在 §3 中回顾。
不过，库仑分支有效场论（EFT）是一个内涵极为丰富的对象：它对库仑分支上支持的
Bogomol'nyi--Prasad--Sommerfield（BPS）粒子谱高度敏感，而当我们穿行模空间时，
% [Tutor Note: 原文此处把 Sommerfeld 拼作 "Sommerfield"（多一个 r），系底本笔误；
%  正确的姓氏为 Arnold Sommerfeld。正文照原样排印，正确形式为 Bogomol'nyi--Prasad--Sommerfeld（BPS）。]
这些 BPS 粒子会经历非平凡的穿墙与单值性 [20,21,35,72]。
% [Tutor Note: 「CB 上远红外是阿贝尔规范理论」是 $\cN=2$ 的强约束：在库仑分支的
%  一般点上，规范群 $G$ 被它的极大环面 $\mathrm{U}(1)^{\rank G}$ 破缺，理论只含无质量
%  阿贝尔矢量多重态；非阿贝尔结构（W 玻色子、磁单极子）都获得质量，作为 BPS 粒子
%  进入有效作用量。这也是 §3 中预势 $\cF(a)$ 只用 $\mathrm{U}(1)$ 的标量 $a^i$ 表述的原因。]
% [原文 p.5]

这种结构在两类可观测量之间有着引人注目的相互作用：
一类自然定义在库仑分支原点处的 SCFT 中——那里 BPS 粒子变得无质量、相互作用变强；
另一类定义在弦/M/F 理论中「工程化」该理论的几何背景里——在那里库仑分支翻译为构造的几何模，
如图 1 所总结。SCFT、其 EFT 与相应的弦/M/F 理论几何之间这种交织的联系，
使得人们可以从场论与几何两种技术去研究 SCFT 的可观测量。
这两种技术之间的对话不仅加深了我们对 SCFT 的物理理解，还导出了数学中许多有趣的关系与恒等式。
例如，库仑分支 EFT 编码了不动点处 BPS 算子的谱，这些算子在共形对称性自发破缺的模空间上
产生 BPS 态 [72--74]。EFT 还通过反常匹配机制 [75] 决定了 SCFT 的各种 't Hooft 反常。
更一般地，EFT 的可观测量与 SCFT 的超对称配分函数密切相关；对拉格朗日量理论，
后者可以通过局域化被精确计算 [76--78]（另见综述 [79] 及其中的文献）。
此外，库仑分支 EFT 提供了 $\cN=2$ SCFT 与几何之间的一种优美互动：
存在一个纤维化在库仑分支之上的涌现黎曼面，由此产生 Seiberg--Witten（SW）几何 [20,21]。
这个黎曼面被称为 Seiberg--Witten 曲线，其复结构在库仑分支上逐点变化，并编码了决定 EFT 的预势。
SW 几何（及其推广）最初被视为求解 EFT 的一种技巧，如今却发现它自然地出现在
由弦/M/F 理论构造 $\cN=2$ 理论的过程中：它记录了（存在膜时）保持 $\cN=2$ 超对称、
并能存活到场论极限的几何背景的某些形变模。§4 的很大一部分将致力于说明：
这样的图像如何从使用 M5 膜（更准确地说，六维 $(2,0)$ 理论）的 Class S 构造中产生；
而在 §5 中，我们将看到一个由 IIB 型弦论探针三重复形奇点所涌现的广义 SW 几何。
在所有这些情形中，场论中的 BPS 解（例如瞬子与粒子）都对应 BPS 膜（与弦）的构型，
而这些构型可以用高维几何的某些枚举不变量来计数 [80--83]。
例如，对一个由 IIB 型弦论探针奇异的 Calabi--Yau（CY）三重复形而工程化的 $\cN=2$ 理论，
其预势由该奇异 CY 的（精化）拓扑弦配分函数决定，而后者对应场论的瞬子配分函数
（关于这一主题的综述见 [84]）。
% [Tutor Note: 「涌现黎曼面」的直观图像是：库仑分支上每个点对应一个物理真空，
%  该真空的电磁耦合与 BPS 谱由一条椭圆曲线的周期给出；因此把库仑分支当作底空间时，
%  每点纤维化一条环面（椭圆曲线），整体就是一个纤维化在 CB 上的黎曼面族。]
% [Tutor Note: 「反常匹配」在这里的含义是：'t Hooft 反常沿重整化群流不变，
%  故可以在 UV（弱耦合、有拉格朗日量）处算出，再「匹配」到 IR 的 SCFT，
%  从而得到 IR 的中心荷 $a,c$ 与味对称性的反常系数 $k_{\mathrm{flavor}}$（见 §2.6 与 §5.1）。]
% [原文 p.6--7]

尽管缺少微扰描述，人们对 $\cN=2$ SCFT 的普适性质在抽象层面上的理解也在稳步推进。
这得益于 SCFT 的公理化定义：由局域算子谱与算符乘积展开（OPE）来定义理论，
而这些数据受结合律、幺正性与超共形对称性的约束。近年来的一个重大进展是
（重新活跃起来的）共形自举纲领（综述见 [85]）：它利用这些约束，
在无穷维理论空间的某个切片上刻画出 CFT 的允许区域。对一般 CFT，最强的界来自
对自举方程的数值求解。而在存在 $\cN=2$ 超共形对称性时，四维 SCFT 含有一个可解而又丰富的
算子子扇区，它在 OPE 下封闭，并由一个涌现的二维手征代数 [86] 描述。
尝试理解与 $\cN=2$ SCFT 相关的这些手征代数空间的「迷你自举」纲领已取得大量进展
（例如见 [87--89]，以及最近的综述 [90]）。此外，给定 SCFT 中手征代数的知识，
还可以把它输入数值自举纲领，从而对理论中更一般的 OPE 数据给出更强的界。
后一种策略同样适用于其他可以从库仑分支 EFT 与上文回顾的弦/M/F 理论几何得到的 SCFT 数据，
尽管在 $\cN=2$ SCFT 中这方面的探索还不算多（见 [91--94]）。
自然地，人们也会追问手征代数在弦/M/F 理论构造中的几何含义。
我们希望，共形自举、EFT 方法与几何构造这三者之间的有力协同，
将大大加深我们对四维 $\cN=2$ SCFT 及其之外的理解。
% [Tutor Note: 由 $\cN=2$ 超共形对称性可以构造出「手征代数」：取特定上同调（Schur 算子）
%  后，四点函数等 OPE 数据退化为一个二维顶点算子代数的数据，且该代数的中心荷由
%  四维中心荷 $c$ 决定（$c_{2d}=-12c_{4d}$）。这就是「solvable sub-sector」的含义：
%  可解，是因为二维手征代数比四维 SCFT 本身容易处理得多。]

% [原文 p.7]
\begin{center}
  \includegraphics[width=0.66\linewidth]{fig1.png}
\end{center}
\figtitle{1}{研究四维 $\cN=2$ SCFT 的各种进路背后的总体思路。图中央是本讲义所讨论的
三条紧密相关而又互补的进路：自举思路、EFT 方法与几何工程。对每条进路，
其基本构件列在邻近的圆圈中；这些构件进一步产生出各自具有自然解释的可观测量。
最外圈列出的是这些可观测量的若干代表，它们描述同一个 SCFT 的不同侧面，
并因此满足非平凡的数学关系。}
% [Tutor Note: 这张图应当这样读：内圈三个节点（Bootstrap / EFT / Geometry）是方法，
%  中圈是每条方法的「输入」（例如 EFT 一侧输入 Moduli space、BPS states、Effective action），
%  外圈是「输出」的可观测量（Anomalies、Protected operators、Operators spectrum、
%  Correlation functions 等）。同一个 SCFT 的所有输出之间由非平凡的恒等式联系，
%  这正是「bootstrap + EFT + geometry 三者互补」的图示表达。]
\clearpage   % 原文 p.8 整页只有图 1，§2 自下一页（PDF p.9）起
